\documentclass[11pt,a4paper]{article}
\usepackage{xeCJK}
\usepackage{fontspec}
\usepackage{xcolor}
\usepackage{fancyhdr}
\usepackage{booktabs}
\usepackage{array}
\usepackage{amssymb}
\usepackage{geometry}
\usepackage[protrusion=true]{microtype}
\geometry{margin=1.8cm,top=2cm,headheight=15pt}
\linespread{1.05}

\setmainfont{Baskerville}
\setCJKmainfont{Songti SC}
\newfontfamily{\starfont}{Songti SC}

\definecolor{wrongred}{RGB}{198,40,40}
\definecolor{wrongbg}{RGB}{253,235,233}
\definecolor{rulegray}{RGB}{200,200,200}

\setlength{\emergencystretch}{3em}
\setlength{\fboxsep}{3pt}
\setlength{\parindent}{0pt}

\newcommand{\checkmarkbox}{\fbox{\rule{0pt}{0.8ex}\rule{0.8ex}{0pt}}}
\newcommand{\STAR}{\textcolor{wrongred}{\starfont ★}}
\newcommand{\bk}{\rule[-0.55ex]{1.0em}{0.4pt}}
\newcommand{\A}[1]{\textbf{#1}}
\newcommand{\W}[1]{\textcolor{wrongred}{#1}}
\newcommand{\OK}{\textcolor{wrongred}{\ensuremath{\checkmark}}}
\newcommand{\NO}{\textcolor{wrongred}{$\times$}}
\newcolumntype{L}[1]{>{\raggedright\arraybackslash}p{#1}}

\pagestyle{fancy}
\fancyhf{}
\fancyhead[L]{\footnotesize 英语语法 动名词练习填空（with-keys-1 · 订正）}
\fancyhead[R]{\footnotesize 赵文瀚}
\fancyfoot[C]{\thepage}
\renewcommand{\headrulewidth}{0.4pt}
\renewcommand{\footrulewidth}{0pt}

\begin{document}

\begin{center}
{\LARGE\bfseries 动名词练习填空（with-keys-1）}\\[0.25em]
{\large 错题订正与考点复盘}\\[0.35em]
{\footnotesize 2029 届高一英语 · 上海市大同中学 · 非谓语动词专项（共 50 题）· 2026-10}
\end{center}

\vspace{-0.2em}
\noindent{\color{rulegray}\rule{\textwidth}{0.8pt}}

\vspace{0.5em}
\noindent\fcolorbox{black!15}{black!4}{%
  \parbox{\dimexpr\linewidth-2\fboxsep-2\fboxrule\relax}{%
    \small 日期：2026-10\qquad 来源：上海市大同中学高一英语「动名词练习填空-with-keys-1」（作业号 260642，共 50 题）\qquad 姓名：赵文瀚\\[0.3em]
    判错：\textcolor{wrongred}{\textbf{23}}\,/\enspace 50\qquad
    判对：\textbf{27}\,/\enspace 50\qquad
    （原卷右上角红笔 \textbf{58.0\%}）\qquad
    重做日期：\underline{\hspace{2.4cm}}}}

\vspace{0.6em}
\noindent{\footnotesize 说明：原卷已由老师红笔批改，下表「判定」\OK\ 判对、\NO\ 判错，\textbf{照录红笔}；\W{红色}＝被判错的作答，\A{加粗}＝正确答案。原卷虽名 with-keys 但\textbf{未印答案}，「正确答案」除红笔订正外为按标准语法补拟，存疑处见「原卷旁注」。第 14、38 题一句两空，判定栏按空分别记。}

\vspace{0.4em}
\renewcommand{\arraystretch}{1.0}
\noindent{\footnotesize\begin{tabular}{@{}L{0.75cm}@{\hspace{0.08cm}}L{8.55cm}@{\hspace{0.08cm}}L{2.25cm}@{\hspace{0.08cm}}L{2.45cm}@{\hspace{0.08cm}}L{0.72cm}@{\hspace{0.08cm}}L{2.25cm}@{}}
\toprule
\textbf{题号} & \textbf{题干（含空与提示词）} & \textbf{我的作答} & \textbf{正确答案} & \textbf{判定} & \textbf{错因} \\
\midrule
1  & Lights are used to \bk\ (control) traffic in cities. & control & control & \OK & — \\
2  & It's difficult for foreigners to get used to \bk\ (eat) with chopsticks. & eating & eating & \OK & — \\
3  & Would you like \bk\ (clean) the room with us? There is a lot of work in it. & to clean & to clean & \OK & — \\
4  & He succeeded in \bk\ (achieve) his aim at last. & achieving & achieving & \OK & — \\
5  & Ryan always attempts to escape \bk\ (punish) when he breaks traffic rules. & \W{punishing} & \A{being punished} & \NO & 动名词被动（escape being done） \\
6  & They don't allow \bk\ (smoke) in this theatre. & smoking（涂改） & smoking & \OK & — \\
7  & They aren't allowed \bk\ (take) the books out of the reading room. & to take & to take & \OK & — \\
8  & I'm sorry I can't help you because I have a lot of letters \bk\ (answer). & to answer & to answer & \OK & — \\
9  & Such kind of person is difficult \bk\ (deal) with. So you'd better leave them alone. & to deal & to deal & \OK & — \\
10 & Except for the old temples, the place is worth \bk\ (see). & seeing & seeing & \OK & — \\
11 & This novel is not worthy of \bk\ (make) into a film. & \W{making} & \A{being made} & \NO & 动名词被动（worthy of being done） \\
12 & The farmer workers were made \bk\ (do) too much without pay. & to do & to do & \OK & — \\
13 & The patient insisted on \bk\ (operate) on even if there is little hope of success. & \W{operating} & \A{being operated} & \NO & 动名词被动（主语是受动者） \\
14 & We failed \bk\ (get) him \bk\ (come) to the party. & to get; \W{coming} & to get; \A{to come} & \OK\,/\NO & get sb to do（第二空错） \\
15 & The doctor has done everything he could \bk\ (save) the wounded soldier. & to save & to save & \OK & — \\
16 & Do you find it any good \bk\ (quarrel) with him about the matter? & \W{quarrelling} & \A{quarrelling} & \NO & 照录红笔；存疑（见旁注） \\
17 & I'd rather \bk\ (not go) out. It's cold outside & not go & not go & \OK & — \\
18 & He preferred \bk\ (die) rather than \bk\ (give) in. & \W{dying; giving} & \A{to die; give} & \NO & prefer to do rather than (to) do \\
19 & Mrs. Frost was worried. I regretted \bk\ (tell) her that her sonwas in the police station & telling & telling & \OK & — \\
20 & I regret \bk\ (tell) you that your son is in the police station. I wonder what has happened to him. & to tell & to tell & \OK & — \\
21 & The little boy did nothing but \bk\ (watch) TV the whole day. & \W{watching} & \A{watch} & \NO & do nothing but do \\
22 & The little boy had no choice but \bk\ (do) what the bigger boys had told him to. & to do & to do & \OK & — \\
23 & Be careful while crossing the road. Don't forget \bk\ (knock) over by a car last month & \W{to knock} & \A{having been knocked} & \NO & forget having been done（已发生） \\
24 & What he wants is \bk\ (not interrupt) while he is speaking. & not to be interrupted & not to be interrupted & \OK & — \\
25 & Because of lack of rain, many people in the southern part of the country have been forced \bk\ (leave) their homes. & to leave & to leave & \OK & — \\
26 & He is devoted to \bk\ (set) up more schools for poor children. & setting & setting & \OK & — \\
\bottomrule
\end{tabular}}

\newpage
\renewcommand{\arraystretch}{1.0}
\noindent{\footnotesize\begin{tabular}{@{}L{0.75cm}@{\hspace{0.08cm}}L{8.55cm}@{\hspace{0.08cm}}L{2.25cm}@{\hspace{0.08cm}}L{2.45cm}@{\hspace{0.08cm}}L{0.72cm}@{\hspace{0.08cm}}L{2.25cm}@{}}
\toprule
\textbf{题号} & \textbf{题干（含空与提示词）} & \textbf{我的作答} & \textbf{正确答案} & \textbf{判定} & \textbf{错因} \\
\midrule
27 & This time my goal is \bk\ (climb) to the mountain top within one hour. & to climb & to climb & \OK & — \\
28 & Her hobby is \bk\ (design) clothes for her dolls, which may seem a little strange to some adults. & \W{designning} & \A{designing} & \NO & 拼写（双写 n） \\
29 & The meeting \bk\ (hold) tomorrow is of great importance & \W{holding} & \A{to be held} & \NO & 不定式被动作后置定语（将来） \\
30 & Tom is said \bk\ (translate) an English novel into Chinese these days. & \W{to translate} & \A{to be translating} & \NO & 不定式进行式（these days） \\
31 & How about the two of us \bk\ (take) a walk after dinner? & taking & taking & \OK & — \\
32 & One needs a lot of practice \bk\ (speak) English well. & to speak（涂改） & to speak & \OK & — \\
33 & These pages need \bk\ (check) again. There might still be some mistakes. & \W{checking}（原写 to be checked） & \A{checking}（to be checked 亦可） & \NO & need doing 主动表被动 \\
34 & He needs \bk\ (practise) typing at least two hours a day. & to practise & to practise & \OK & — \\
35 & Jane happened \bk\ (be) there before, so we might as well ask herto be our guide. & \W{to be} & \A{to have been} & \NO & 不定式完成式（先于 happened） \\
36 & I am so glad \bk\ (give) a chance to visit your beautiful city. & \W{giving}（原写 to be given） & \A{to have been given} & \NO & 不定式被动完成；划改反改错 \\
37 & Bob was seen \bk\ (fall) from the tree the other day. & to fall & to fall & \OK & — \\
38 & It is no use \bk\ (hide) it from your mother. She appears \bk\ (tell) everything. & hiding; \W{to tell} & hiding; \A{to have been told} & \OK\,/\NO & appears to have been done \\
39 & Brad pretended \bk\ (work) on his term paper when mother came in. & \W{to work} & \A{to be working} & \NO & 不定式进行式（when 从句背景） \\
40 & Try \bk\ (add) some vinegar to the soup. Maybe it will taste better. & adding & adding & \OK & — \\
41 & He blamed me for \bk\ (not write) to him for nearly a year. & \W{not writing} & \A{not having written} & \NO & 动名词完成式（先于 blamed） \\
42 & My suggestion couldn't help \bk\ (improve) their service. & to improve & to improve & \OK & — \\
43 & Seeing a cafe by the roadside, he stopped \bk\ (have) a cup of coffee there. & to have & to have & \OK & — \\
44 & She hurried to the station only \bk\ (tell) that the train had left & \W{telling}（原写 was told） & \A{to be told} & \NO & only to do 意外结果（被动） \\
45 & I suggested my brother \bk\ (run) a business of his own & \W{run} & \A{(should) run} & \NO & suggest 虚拟语气；存疑（见旁注） \\
46 & The goods are believed \bk\ (ship) secretly to another country a few days ago. & \W{shipping} & \A{to have been shipped} & \NO & 不定式被动完成（a few days ago） \\
47 & We thought it wrong for her \bk\ (punish) & \W{punishing}（原写 to be punished） & \A{to be punished} & \NO & for sb to be done；划改反改错 \\
48 & \bk\ (admit) into the Party, you have to write an application first. & \W{Being admitted} & \A{To be admitted} & \NO & 不定式被动表将来/条件 \\
49 & I will never forgive \bk\ (he insult) me in public & his insulting & his insulting & \OK & — \\
50 & Something has been done to prevent workers from \bk\ (injure) when they are working & \W{injuring}（原写 being injured） & \A{being injured} & \NO & 介词后动名词被动；划改反改错 \\
\bottomrule
\end{tabular}}

\vspace{0.3em}
\noindent{\footnotesize\textbf{原卷旁注}：（一）右上角红笔「58.0\%」为作业系统记录；按红笔逐题清点为 \OK\,27、\NO\,23（54\%），且第 14、38 题一句两空各对一空，与 58.0\% 口径略有出入，请以老师讲评为准。（二）第 6、32 题作答有涂改（第 6 题先写 \emph{the} 后涂掉，第 32 题先写别词后涂掉补 \emph{to speak}），红笔均判对。（三）第 33、36、44、47、50 题学生自行划改后再答，红笔按最终作答判错；其中第 36、47 题最初的写法（\emph{to be given}／\emph{to be punished}）恰是正确形式。（四）第 16 题 \emph{quarrelling}、第 41 题 \emph{not writing}、第 45 题 \emph{run} 被红笔判错，但作答形似常见答案：第 41 题应取 \emph{not having written}（先于 \emph{blamed}），第 45 题或按 \emph{suggest}＋(should) do／\emph{my brother's running} 判，请对照讲评核实。（五）原卷印题笔误照录：第 12 题 \emph{The farmer workers}、第 19 题 \emph{sonwas}、第 35 题 \emph{herto}、第 23 题 \emph{(knock)over}。（六）第 12 题 \emph{were made to do}（被动式 make 后接 to 不定式）红笔判对；第 42 题 \emph{help to do}（有助于）红笔判对。}

\newpage
\begin{center}
{\Large\bfseries 错因归类与复盘}
\end{center}
\vspace{-0.4em}
\noindent{\color{rulegray}\rule{\textwidth}{0.8pt}}

\vspace{0.5em}
\noindent 说明：按非谓语动词考点归类 23 条错题（第 14、38 题一句两空，各按错空计一处），供汇总对账；逐题判定见第 1--2 页。

\vspace{0.5em}
\renewcommand{\arraystretch}{1.35}
\begin{center}
\begin{tabular}{@{}p{6.6cm} p{7.6cm} c@{}}
\toprule
\textbf{易错考点} & \textbf{题号} & \textbf{题数} \\
\midrule
不定式的被动／进行／完成式（to be done／to be doing／to have (been) done） & {\small 23, 29, 30, 35, 36, 38, 39, 44, 46, 47, 48} & 11 \\
动名词的被动／完成式（being done／having (been) done） & {\small 5, 11, 13, 33, 41, 50} & 6 \\
固定搭配（but do／rather than (to) do／get sb to do／suggest＋虚拟） & {\small 14, 18, 21, 45} & 4 \\
词形拼写（designing） & {\small 28} & 1 \\
照录判错、作答存疑（quarrelling） & {\small 16} & 1 \\
\midrule
\textbf{合计} & \textbf{23 题} & \textbf{23} \\
\bottomrule
\end{tabular}
\end{center}

\vspace{0.5em}
\noindent\textbf{关联知识点标签}
\vspace{0.25em}

\noindent\begin{tabular}{@{}p{8.2cm}@{\hspace{0.4cm}}p{8.2cm}@{}}
\begin{minipage}[t]{\linewidth}
\small
\STAR\ \textbf{动名词被动式}：escape／insist on／worthy of＋\textbf{being done}（主语是受动者）；need \textbf{doing}（＝need to be done，主动表被动）；介词 from＋\textbf{being done}\\[0.2em]
\STAR\ \textbf{动名词完成式}：regret／blame sb for＋\textbf{having done}（动作先于主句）\\[0.2em]
\STAR\ \textbf{need 两用}：need \textbf{doing}／need \textbf{to be done}（本卷老师按前者判）；sb need \textbf{to do}（人作主语）
\end{minipage} &
\begin{minipage}[t]{\linewidth}
\small
\STAR\ \textbf{不定式三式}：is said \textbf{to be doing}（进行）；happen／appear \textbf{to have (been) done}（先于主句）；only \textbf{to be done}（意外结果）\\[0.2em]
\STAR\ \textbf{固定搭配}：do nothing but \textbf{do}／have no choice but \textbf{to do}；prefer \textbf{to do} rather than \textbf{(to) do}；get sb \textbf{to do}；suggest＋(should) \textbf{do}（虚拟）\\[0.2em]
\STAR\ \textbf{被动式 make／see}：be made \textbf{to do}；be seen \textbf{to do}（被动后还原 to）
\end{minipage} \\
\end{tabular}

\vspace{0.65em}
\noindent\fcolorbox{black!15}{black!4}{%
  \parbox{\dimexpr\linewidth-2\fboxsep-2\fboxrule\relax}{%
    \textbf{复盘记录}\\[0.35em]
    最薄弱的板块：\underline{\hspace{9.0cm}}\\[0.55em]
    主要错误原因：\checkmarkbox 被动式漏 being/to be\quad \checkmarkbox 完成式漏 having/to have\quad \checkmarkbox 固定搭配\quad \checkmarkbox 拼写\\[0.55em]
    本次复习的改进措施：\underline{\hspace{10.1cm}}\\[0.55em]
    \hspace{3.2em}\underline{\hspace{10.1cm}}\\[0.55em]
    下次复习日期：\underline{\hspace{2.2cm}}\qquad
    当前掌握度：\checkmarkbox 仍薄弱\quad \checkmarkbox 基本掌握\quad \checkmarkbox 已掌握
  }}

\end{document}
